\phantomsection
\addcontentsline{toc}{part}{第二部分\hspace{1em}文献综述和开题报告}
\phantomsection
\addcontentsline{toc}{sectionnoiddot}{指导教师对文献综述和开题报告具体内容要求}
\noindent{\setfont{4}{1.5}\selectfont\bfseries {一、题目： \ProjectName}} \\[14bp]
\noindent{\setfont{4}{1.5}\selectfont\bfseries {二、指导老师对文献综述、开题报告、外文翻译的具体要求}}
\vspace{7bp}
\par
% ------ 不要动上面的 ------

文献综述要求：围绕多普勒雷达手势识别主题，重点覆盖（1）连续波多普勒雷达的信号模型与近场散射理论；（2）基于视觉、可穿戴及雷达的手势识别方法综述；（3）机器学习分类方法（SVM、CNN、RNN/LSTM）在手势识别中的应用现状与对比。字数不少于3000字，参考文献不少于15篇，其中英文文献不少于10篇。

开题报告要求：在文献综述基础上，明确本课题的研究目标、技术路线与可行性分析，说明已开展的阶段性工作（包括硬件搭建与数据采集初步结果），给出详细的进度安排与预期目标。

外文翻译要求：翻译指定参考论文（Y. Zhang et al., IEEE TMT 2021），全文翻译，字数不少于5000汉字，附原文。

% ------ 不要动下面的（只改日期即可） ------
\newcommand{\foottext}[1]{{\setfont{-4}{1}\selectfont\bfseries {#1}}}
\vfill
\hfill \makebox[\widthof{\foottext{指导教师（签名）}}][l]{\foottext{指导教师（签名）}}\underline{\makebox[3cm][c]{\foottext{ }}}\par
\vspace{6bp}
\hfill \foottext{2025年11月25日}
\let\foottext\undefined